% Collocamento: il regolamento in formato A5, per il libretto del print and play.
% Le pagine vanno a gruppi di quattro: grafica/pnp.py le impagina due per foglio A4,
% fronte e retro, da piegare a metà.
%
% Le immagini e i nomi dei poteri li prepara figure.py (py -3 pitch/figure.py), poi:
%   pdflatex regolamento_a5.tex   (due volte)
\documentclass[9pt]{extarticle}
\usepackage[T1]{fontenc}
\usepackage[utf8]{inputenc}
\usepackage[italian]{babel}
\usepackage{tgheros}
\renewcommand{\familydefault}{\sfdefault}
\usepackage[a5paper,left=19mm,right=11mm,top=19mm,bottom=17mm]{geometry}
\usepackage{xcolor,graphicx,tikz,enumitem,array,tabularx,colortbl,eso-pic,calc,pgffor}
\usepackage[most]{tcolorbox}
\usepackage[tracking=true,expansion=false]{microtype}
\usetikzlibrary{calc,arrows.meta}
% Scritto da pitch/figure.py a partire da grafica/dati/poteri.json: non modificarlo a mano.
\newcommand{\potereUno}{Passaparola}
\newcommand{\effettoUno}{Pesca 1 carta Ambito.}
\newcommand{\breveUno}{pesca 1 Ambito.}
\newcommand{\chiUno}{08 · 10 · 12 · 13 · 15 · 18 · 23 · 24 · 33 · 34 · 35 · 49 · 55 · 57 · 58}
\newcommand{\potereDue}{Colloquio}
\newcommand{\effettoDue}{Prendi una Persona dal mercato e tienila da parte, scoperta: è prenotata per te. Il mercato si ripristina.}
\newcommand{\breveDue}{prenota una Persona del mercato.}
\newcommand{\chiDue}{04 · 07 · 09 · 14 · 19 · 28 · 38 · 42 · 44 · 46 · 47 · 52 · 53 · 54 · 60}
\newcommand{\potereTre}{Corso intensivo}
\newcommand{\effettoTre}{Gioca subito fino a 2 carte Ambito dalla mano come Formazioni: per aprire pile nuove o per migliorarne di tue.}
\newcommand{\breveTre}{gioca fino a 2 Ambiti come formazioni.}
\newcommand{\chiTre}{05 · 16 · 17 · 25 · 26 · 32 · 37 · 39 · 43 · 45 · 48 · 50 · 51 · 56 · 59}
\newcommand{\potereQuattro}{Chiamata diretta}
\newcommand{\effettoQuattro}{Metti subito un'altra Persona, dal mercato o tra quelle prenotate, su una tua pila libera, ignorando una delle sue tre ambizioni. Il suo potere non si attiva.}
\newcommand{\breveQuattro}{piazza un'altra Persona ignorando un'ambizione (senza il suo potere).}
\newcommand{\chiQuattro}{01 · 02 · 03 · 06 · 11 · 20 · 21 · 22 · 27 · 29 · 30 · 31 · 36 · 40 · 41}


\pdfinfo{/Title (Collocamento - Regolamento)}

\definecolor{carta}{HTML}{F3ECDD}
\definecolor{cartascura}{HTML}{E5DACA}
\definecolor{petrolio}{HTML}{17474D}
\definecolor{rosa}{HTML}{FF3D8B}
\definecolor{magenta}{HTML}{D53883}
\colorlet{accento}{rosa}
\tcbset{colupper=petrolio}

% ---------------------------------------------------------------- la pagina è una carta
\newif\ifpaginacarta \paginacartatrue
\newcommand{\duecifre}[1]{\ifnum#1<10 0\fi\number#1}
\AddToShipoutPictureBG{\AtPageLowerLeft{%
  \ifpaginacarta
  \begin{tikzpicture}[x=1mm,y=1mm]
    \useasboundingbox (0,0) rectangle (148,210);
    \fill[carta] (0,0) rectangle (148,210);
    \fill[cartascura] (0,0) rectangle (10,210);
    \draw[petrolio,line width=1.2pt] (10,0) -- (10,210);
    \node[rotate=90,text=petrolio,font=\fontsize{5}{6}\selectfont,inner sep=0] at (5,80) {\textls[300]{COLLOCAMENTO · REGOLAMENTO}};
    \node[text=petrolio,font=\fontsize{4.6}{5}\selectfont\bfseries,inner sep=0] at (5,8.4) {CDS};
    \node[anchor=base west,text=petrolio,font=\fontsize{6}{7}\selectfont\bfseries,inner sep=0] at (19,199.6) {\textls[240]{COLLOCAMENTO}};
    \node[anchor=base east,text=petrolio,font=\fontsize{6.8}{7}\selectfont,inner sep=0] at (137,199.6) {\duecifre{\value{page}}};
    \draw[petrolio,line width=.8pt] (19,197) -- (137,197);
    \draw[petrolio,line width=.3pt] (19,11.6) -- (137,11.6);
    \node[anchor=base west,text=petrolio,font=\fontsize{4.6}{5}\selectfont\bfseries,inner sep=0] at (19,8.2) {CDS · COLLOCAMENTO};
    \node[anchor=base east,text=petrolio,font=\fontsize{4.6}{5}\selectfont,inner sep=0] at (137,8.2) {REGOLAMENTO · PROTOTIPO · 2026};
  \end{tikzpicture}%
  \fi}}

\pagestyle{empty}
\color{petrolio}
\setlength{\parindent}{0pt}
\setlength{\parskip}{1.5mm}
\linespread{1.06}
\raggedbottom
\setlist[itemize]{leftmargin=3.8mm,itemsep=.5mm,topsep=.4mm,parsep=0pt,partopsep=0pt,label={\color{accento}\rule[.22ex]{1.2mm}{1.2mm}}}
\renewcommand{\arraystretch}{1.35}

\newcommand{\etichetta}[1]{{\fontsize{5.6}{7}\selectfont\bfseries\textls[170]{\MakeUppercase{#1}}}}
\newcommand{\titolone}[1]{\par\noindent{\color{accento}\rule{10mm}{2.2pt}}\par\vspace{1mm}\noindent{\fontsize{16}{18}\selectfont\bfseries\MakeUppercase{#1}}\par\vspace{2mm}}
\newcommand{\titolo}[1]{\par\vspace{3mm}\noindent{\color{accento}\rule{8mm}{1.9pt}}\par\vspace{.6mm}\noindent{\fontsize{11.5}{13}\selectfont\bfseries\MakeUppercase{#1}}\par\vspace{.8mm}}
\newcommand{\sottotitolo}[1]{\par\vspace{1.6mm}\noindent{\fontsize{9}{10.5}\selectfont\bfseries #1}\par\vspace{-.6mm}}
\newcommand{\attacco}[1]{{\fontsize{10}{13.5}\selectfont #1\par}}

\newtcolorbox{nota}[1]{enhanced,colback=cartascura,colframe=cartascura,boxrule=0pt,arc=0pt,outer arc=0pt,
  left=3.4mm,right=2.8mm,top=2mm,bottom=2mm,boxsep=0pt,borderline west={2pt}{0pt}{accento},
  before skip=2.4mm,after skip=2.4mm,parbox=false,
  before upper={{\color{accento}\etichetta{#1}}\par\vspace{.2mm}}}

\newcounter{passo}
\newcommand{\daccapo}{\setcounter{passo}{0}}
\newcommand{\passo}[2]{\stepcounter{passo}\par\vspace{.3mm}\noindent
  \begin{minipage}[t]{5.6mm}\vspace{0pt}\tikz\node[fill=petrolio,text=white,font=\fontsize{8.5}{8.5}\selectfont\bfseries,minimum size=5mm,inner sep=0]{\thepasso};\end{minipage}\hspace{2mm}%
  \begin{minipage}[t]{\dimexpr\linewidth-7.6mm\relax}\vspace{.1mm}\textbf{#1}\enspace #2\end{minipage}\par\vspace{.5mm}}

\newcommand{\ico}[1]{\raisebox{-.26em}{\includegraphics[height=1.2em]{img/ico_#1}}}
\newcommand{\terna}[1]{\foreach \s in {#1}{\ico{\s}\hspace{.1em}}}
\newcommand{\spunta}{\tikz[baseline=.1mm]\draw[line width=1pt,accento,line cap=round,line join=round] (0,.9mm) -- (.75mm,0) -- (2.2mm,2.3mm);}
\newcommand{\senzasillabe}{\hyphenpenalty=10000 \exhyphenpenalty=10000 }

\newcommand{\cartaAt}[5][]{%
  \begin{scope}[shift={(#2)},rotate=#3]
    \pgfmathsetmacro{\cartaW}{#4/2}\pgfmathsetmacro{\cartaH}{#4*88/126}\pgfmathsetmacro{\cartaR}{#4*0.04}
    \fill[black,opacity=.22,rounded corners=\cartaR mm] (-\cartaW+.45,-\cartaH-.6) rectangle (\cartaW+.45,\cartaH-.6);
    \begin{scope}
      \clip[rounded corners=\cartaR mm] (-\cartaW,-\cartaH) rectangle (\cartaW,\cartaH);
      \node[inner sep=0,outer sep=0,rotate=#3] at (0,0) {\includegraphics[width=#4mm,height=\dimexpr#4mm*88/63\relax]{img/#5}};
    \end{scope}
    \draw[petrolio,opacity=.6,line width=.25pt,rounded corners=\cartaR mm,#1] (-\cartaW,-\cartaH) rectangle (\cartaW,\cartaH);
  \end{scope}}
\newcommand{\sullacarta}[6]{\coordinate (#1) at ({#2+(#5/756-.5)*#4},{#3+(.5-#6/1056)*#4*88/63});}

\begin{document}

% ================================================================ COPERTINA
\paginacartafalse
\AddToShipoutPictureBG*{\AtPageLowerLeft{\begin{tikzpicture}[x=1mm,y=1mm]
  \useasboundingbox (0,0) rectangle (148,210);
  \fill[carta] (0,0) rectangle (148,210);
  \fill[rosa] (0,0) rectangle (148,62);
  \draw[petrolio,line width=1.6pt] (0,62) -- (148,62);
  \node[text=petrolio,font=\fontsize{8}{10}\selectfont,inner sep=0] at (74,186) {\textls[480]{CDS}};
  \node[text=petrolio,font=\fontsize{29}{31}\selectfont\bfseries,inner sep=0] at (74,170) {COLLOCAMENTO};
  \draw[accento,line width=1.4pt] (44,159) -- (104,159);
  \node[text=petrolio,font=\fontsize{9}{11}\selectfont,inner sep=0] at (74,152) {\textls[200]{REGOLAMENTO}};
  \node[text=petrolio,font=\fontsize{8.6}{12}\selectfont\itshape,align=center,text width=110mm,inner sep=0] at (74,139) {Forma le persone rispettando le loro ambizioni,\\trova a ciascuna il lavoro giusto e chiudi il collocamento\\prima del resto del tavolo.};
  \cartaAt{30,76}{12}{38}{lav_03}
  \cartaAt{118,76}{-12}{38}{lav_02}
  \cartaAt{52,84}{6}{40}{amb_27}
  \cartaAt{96,84}{-6}{40}{amb_17}
  \cartaAt{74,87}{0}{42}{lav_07}
  \node[text=petrolio,font=\fontsize{7.4}{9}\selectfont\bfseries,inner sep=0] at (74,22) {\textls[200]{2-4 PARTECIPANTI · CARTE · A TURNI}};
\end{tikzpicture}}}
\null\clearpage\paginacartatrue

% ================================================================ IL GIOCO IN BREVE
\titolone{Il gioco in breve}

\attacco{Ogni partecipante gestisce un piccolo ufficio di collocamento.}

Con le carte Ambito costruisce percorsi di \textbf{Formazione}. Quando un percorso offre ciò che una \textbf{Persona} desidera, la accoglie. Poi le trova un \textbf{Lavoro} che quel percorso rende possibile. Formazione, Persona e Lavoro formano un \textbf{set}: più il Lavoro somiglia alle ambizioni della Persona, più punti vale.

La fine della partita scatta al primo quinto set chiuso al tavolo. Vince chi ha più punti.

\titolo{Materiale}

\begin{itemize}
  \item \textbf{60 carte Persona.} Ognuna ha tre icone \textbf{ambizione} e un \textbf{potere}, scritto sulla carta.
  \item \textbf{60 carte Ambito:} 30 ambiti, due copie di ognuno. Ogni carta ha tre icone \textbf{competenza} e appartiene a una di cinque \textbf{categorie}: Tecnica e produzione, Servizi e relazioni, Organizzazione e gestione, Creatività e comunicazione, Ricerca e innovazione.
  \item \textbf{4 token obiettivo:} due da 5 punti e due da 3.
  \item \textbf{4 carte riassunto}, una a testa.
\end{itemize}

\clearpage
% ================================================================ COMPETENZE E CARTE
\titolo{Le competenze}

Sono sei, e sono il linguaggio del gioco: ogni Persona e ogni Ambito ne mostra tre. La stessa icona può comparire più volte sulla stessa carta.

\vspace{1mm}
\newcommand{\competenza}[3]{\begin{minipage}[c]{.32\linewidth}\raisebox{-.45\height}{\includegraphics[width=7.5mm]{img/ico_#1}}\hspace{1.6mm}\begin{minipage}[c]{.62\linewidth}{\fontsize{8}{9}\selectfont\bfseries #2\par}{\fontsize{6.4}{8}\selectfont #3\par}\end{minipage}\end{minipage}}
\noindent\competenza{MA}{Manualità}{comune}\hfill\competenza{DI}{Digitale}{comune}\hfill\competenza{CO}{Comunicazione}{comune}

\vspace{2.2mm}
\noindent\competenza{OR}{Organizzazione}{non comune}\hfill\competenza{CR}{Creatività}{non comune}\hfill\competenza{RI}{Ricerca}{rara}

\begin{nota}{Una carta, due usi}
Una carta Ambito si gioca come \textbf{Formazione} oppure come \textbf{Lavoro}. Come Formazione \emph{offre} le sue tre competenze; come Lavoro le \emph{richiede}.
\end{nota}

\titolo{Le carte}

\vspace{1mm}
\noindent\begin{tikzpicture}[x=1mm,y=1mm]
  \useasboundingbox (0,-6) rectangle (118,60);
  \cartaAt{38,30}{0}{38}{lav_17}
  \cartaAt{99,30}{0}{38}{amb_17}
  \begin{scope}[font=\fontsize{6.2}{7.8}\selectfont\senzasillabe,line width=.4pt]
    \sullacarta{la}{38}{30}{38}{0}{36}\sullacarta{lb}{38}{30}{38}{0}{446}
    \draw ($(la)+(-.8,0)$) -- ($(la)+(-1.8,0)$) -- ($(lb)+(-1.8,0)$) -- ($(lb)+(-.8,0)$);
    \node[anchor=east,align=right,text width=16mm,inner sep=0] at ($(la)!.5!(lb)+(-2.6,0)$) {\textbf{Tre ambizioni.}\\Quello che la formazione deve offrirle.};
    \sullacarta{ln}{38}{30}{38}{0}{748}
    \draw ($(ln)+(6.3,0)$) -- ($(ln)+(-1.8,0)$); \fill ($(ln)+(6.3,0)$) circle (.45);
    \node[anchor=east,align=right,text width=16mm,inner sep=0] at ($(ln)+(-2.6,0)$) {\textbf{Nome e battuta.}};
    \sullacarta{lp}{38}{30}{38}{600}{1000}
    \draw ($(lp)+(0,-.8)$) -- ($(lp)+(0,-6.6)$); \fill ($(lp)+(0,-.8)$) circle (.45);
    \node[anchor=north,inner sep=0] at ($(lp)+(0,-7.3)$) {\textbf{Il suo potere}};
    \sullacarta{aa}{99}{30}{38}{0}{36}\sullacarta{ab}{99}{30}{38}{0}{446}
    \draw ($(aa)+(-.8,0)$) -- ($(aa)+(-1.8,0)$) -- ($(ab)+(-1.8,0)$) -- ($(ab)+(-.8,0)$);
    \node[anchor=east,align=right,text width=17mm,inner sep=0] at ($(aa)!.5!(ab)+(-2.6,0)$) {\textbf{Tre competenze.}\\Offerte come Formazione, richieste come Lavoro.};
    \sullacarta{ac}{99}{30}{38}{0}{962}
    \draw ($(ac)+(3,0)$) -- ($(ac)+(-1.8,0)$); \fill ($(ac)+(3,0)$) circle (.45);
    \node[anchor=east,align=right,text width=17mm,inner sep=0] at ($(ac)+(-2.6,0)$) {\textbf{Categoria.}\\Serve per i token.};
  \end{scope}
\end{tikzpicture}


\clearpage
\titolo{Preparazione}
\daccapo
\passo{Mescolate}{separatamente il mazzo delle Persone e quello degli Ambiti.}
\passo{Pescate}{3 carte Ambito a testa: è la vostra mano.}
\passo{Aprite il mercato.}{Rivelate in fila una Persona per ogni partecipante. Le nuove entreranno sempre in coda alla fila.}
\passo{Preparate i token}{accanto ai mazzi e lasciate spazio per gli scarti degli Ambiti.}
\passo{Prendete una carta riassunto}{a testa: da un lato il turno, dall'altro punti e poteri.}

\vspace{1mm}
\noindent\begin{tikzpicture}[x=1mm,y=1mm]
  \useasboundingbox (0,-3) rectangle (118,47);
  \begin{scope}[font=\fontsize{5}{6}\selectfont\bfseries]
    \cartaAt{8,30}{0}{13}{dorsoB_amb}
    \node[inner sep=0] at (8,18.6) {\textls[120]{AMBITI}};
    \draw[petrolio,line width=.4pt,dash pattern=on 1mm off .8mm,rounded corners=.5mm] (17.5,20.9) rectangle (30.5,39.1);
    \node[inner sep=0] at (24,18.6) {\textls[120]{SCARTI}};
    \foreach \x/\y/\v in {37/34/5,44.5/34/5,37/26/3,44.5/26/3}{\fill[petrolio] (\x,\y) circle (3.2);\draw[carta,line width=.5pt] (\x,\y) circle (2.6);\node[text=carta,font=\fontsize{6.5}{6.5}\selectfont\bfseries,inner sep=0] at (\x,\y) {\v};}
    \node[inner sep=0] at (40.75,18.6) {\textls[120]{TOKEN}};
    \cartaAt{62,30}{0}{13}{lav_06}
    \cartaAt{77,30}{0}{13}{lav_28}
    \cartaAt{92,30}{0}{13}{lav_54}
    \cartaAt{110,30}{0}{13}{dorsoB_lav}
    \node[inner sep=0] at (77,42.8) {\textls[120]{IL MERCATO}};
    \node[inner sep=0] at (110,18.6) {\textls[120]{PERSONE}};
    \draw[accento,line width=.8pt,-{Stealth[length=1.8mm]}] (100,15.4) -- (55,15.4);
    \node[inner sep=0,font=\fontsize{5}{6}\selectfont] at (77.5,13) {le nuove entrano a destra, quella in fila da più tempo esce a sinistra};
    \cartaAt{8,4}{10}{11}{amb_13}
    \cartaAt{16,5}{0}{11}{amb_08}
    \cartaAt{24,4}{-10}{11}{amb_17}
    \node[anchor=west,inner sep=0] at (33,4) {\textls[120]{LA TUA MANO: 3 AMBITI}};
  \end{scope}
\end{tikzpicture}

Inizia chi è più giovane, poi si procede in senso orario. Il limite di mano è \textbf{6 carte}.

\clearpage
% ================================================================ IL TURNO
\titolo{Il turno}

Nel tuo turno fai \textbf{una sola} di queste cinque azioni. I poteri delle Persone possono concederne altre.
\daccapo
\passo{Aprire una formazione.}{Gioca una carta Ambito davanti a te: diventa una Formazione e apre una nuova pila. Puoi avere quante pile vuoi.}
\passo{Migliorare una formazione.}{Gioca una carta Ambito sotto una tua Formazione: le sue tre competenze si sommano a quelle della pila. Puoi farlo anche dopo aver inserito una Persona, finché il set non è chiuso.}
\passo{Inserire una persona.}{Prendi una Persona dal mercato e mettila sopra una tua pila che non ne ha ancora una. Le competenze della pila devono comprendere \textbf{tutte e tre le sue ambizioni}, contate una per una: due icone Digitale ne chiedono due. Il mercato si ripristina subito, in coda. Poi \textbf{attivi il potere} della Persona che hai preso.}
\passo{Completare un set con un lavoro.}{Gioca dalla mano una carta Ambito, come Lavoro, sopra una tua pila che ha già una Persona. Le Formazioni della pila devono comprendere \textbf{tutte e tre le icone del Lavoro}. Segna i punti, controlla i token obiettivo e metti negli scarti le Formazioni: Persona e Lavoro restano davanti a te.}
\passo{Pescare.}{Pesca 2 carte Ambito. Poi il mercato scorre: la Persona in fila da più tempo va in fondo al mazzo e ne entra una nuova in coda.}

Se non puoi fare nessuna azione, passi.

\begin{nota}{Fine del turno}
Se hai più di 6 carte in mano, scarta quelle in eccesso. È l'ultima cosa che fai nel turno.
\end{nota}

\clearpage
\titolo{Come si costruisce un set}

Le carte si sovrappongono lasciando visibili tutte le fasce delle icone. Ogni nuova Formazione va \textbf{sotto} le precedenti; la Persona va sopra le Formazioni; il Lavoro va sopra la Persona. In una pila le competenze in più, o ripetute, non danno fastidio.

\vspace{2mm}
\noindent\begin{tikzpicture}[x=1mm,y=1mm]
  \useasboundingbox (0,0) rectangle (118,50);
  \cartaAt{15,22}{0}{30}{amb_08}
  \cartaAt{20.475,22}{0}{30}{amb_13}
  \cartaAt{25.95,22}{0}{30}{lav_17}
  \cartaAt{31.425,22}{0}{30}{amb_17}
  \foreach \x/\n in {2.74/1,8.21/2,13.69/3,19.16/4}{\fill[petrolio] (\x,46.6) circle (1.8);\node[text=white,font=\fontsize{6}{6}\selectfont\bfseries,inner sep=0] at (\x,46.6) {\n};\draw[petrolio,line width=.3pt] (\x,44.8) -- (\x,43.2);}
  \begin{scope}[shift={(50,0)}]
    \newcommand{\riga}[4]{\node[anchor=north west,inner sep=0,text width=68mm,font=\fontsize{7}{9.6}\selectfont] at (0,#1) {{\fontsize{5.6}{7}\selectfont\bfseries\color{accento}\textls[150]{#2}}\\[-.2mm]\textbf{#3}\enspace #4};}
    \riga{49}{1 · 2\enspace FORMAZIONI}{Ristorazione e Logistica.}{La pila offre \terna{MA,CO,CO}\,\terna{MA,DI,OR}}
    \riga{36.5}{3\enspace PERSONA}{Il Tecnico Paziente}{desidera \terna{MA,DI,CO}\,e la pila le ha tutte \spunta}
    \riga{24.5}{4\enspace LAVORO}{Direzione}{richiede \terna{CO,MA,OR}\,e la pila le ha tutte \spunta}
    \riga{12.5}{PUNTI}{Due ambizioni combaciano col Lavoro,}{\ico{MA} e \ico{CO}: il set vale \textbf{5 + 5 = 10 punti}. Con Assistenza clienti \terna{DI,CO,MA} varrebbe 16.}
  \end{scope}
\end{tikzpicture}

\vspace{1mm}
Chiuso il set, le Formazioni vanno negli scarti. Davanti a te restano la Persona e il suo Lavoro: un collocamento riuscito.

\begin{nota}{Quante formazioni servono}
Di solito due. Ne basta una quando Formazione, Persona e Lavoro hanno le stesse tre icone: succede, e vale il massimo dei punti.
\end{nota}

\clearpage
\titolo{Punti del collocamento}

Ogni set vale \textbf{5 punti}, più un bonus per le ambizioni della Persona che combaciano con le icone della carta giocata come Lavoro. Ogni icona si abbina una volta sola.

\vspace{.5mm}
\noindent{\arrayrulecolor{carta}\setlength{\arrayrulewidth}{1pt}%
\begin{tabularx}{\linewidth}{>{\centering\arraybackslash}X|>{\centering\arraybackslash}X|>{\centering\arraybackslash}X|>{\centering\arraybackslash}X}
\rowcolor{petrolio}\textcolor{white}{\etichetta{Corrispondenze}} & \textcolor{white}{\etichetta{Base}} & \textcolor{white}{\etichetta{Bonus}} & \textcolor{white}{\etichetta{Totale}}\\\hline
\rowcolor{cartascura} nessuna & 5 & 0 & \textbf{5}\\\hline
\rowcolor{cartascura} 1 & 5 & 2 & \textbf{7}\\\hline
\rowcolor{cartascura} 2 & 5 & 5 & \textbf{10}\\\hline
\rowcolor{cartascura} 3 & 5 & 11 & \textbf{16}
\end{tabularx}}

\vspace{1.5mm}
Le ambizioni non devono per forza combaciare col Lavoro: un lavoro qualunque si trova sempre. È il lavoro dei sogni che fa la differenza.

% ================================================================ TOKEN E POTERI
\titolo{I token obiettivo}

Per gli obiettivi conta la \textbf{categoria} della carta che ha chiuso ogni set come Lavoro.

\sottotitolo{Tre set della stessa categoria}
Il token da \textbf{5 punti} va a chi chiude, prima del resto del tavolo, 3 set della stessa categoria; quello da \textbf{3 punti} a chi ci riesce subito dopo.

\sottotitolo{Quattro set di categorie diverse}
Il token da \textbf{5 punti} va a chi chiude, prima del resto del tavolo, 4 set di quattro categorie diverse; quello da \textbf{3 punti} a chi ci riesce subito dopo.

\vspace{1mm}
Ogni partecipante può prendere al massimo un token per obiettivo. I token si sommano al punteggio finale.

\clearpage
\titolo{I poteri delle persone}

Ogni Persona ha uno di quattro poteri, scritto sulla carta. Il potere si attiva \textbf{quando inserisci la Persona} con l'azione 3.

\vspace{.5mm}
\noindent{\arrayrulecolor{carta}\setlength{\arrayrulewidth}{1pt}%
\begin{tabularx}{\linewidth}{>{\raggedright\arraybackslash}p{19mm}|>{\raggedright\arraybackslash}X}
\rowcolor{petrolio}\textcolor{white}{\etichetta{Potere}} & \textcolor{white}{\etichetta{Effetto}}\\\hline
\rowcolor{cartascura}\textbf{\potereUno} & \effettoUno\\\hline
\rowcolor{cartascura}\textbf{\potereDue} & \effettoDue\\\hline
\rowcolor{cartascura}\textbf{\potereTre} & \effettoTre\\\hline
\rowcolor{cartascura}\textbf{\potereQuattro} & \effettoQuattro
\end{tabularx}}

\vspace{1.5mm}
\begin{itemize}
  \item \textbf{Una Persona prenotata} con \potereDue{} si inserisce con l'azione 3, al posto di prenderne una dal mercato: valgono gli stessi requisiti e il mercato non cambia. Il suo potere si attiva in quel momento, non quando la prenoti.
  \item Le Persone prenotate non contano nel limite di mano e può usarle solo chi le ha prenotate.
  \item Un potere che non puoi o non vuoi usare va perso.
\end{itemize}

\clearpage
% ================================================================ FINE
\titolo{Fine della partita}

La fine della partita scatta al primo \textbf{quinto set} completato. Terminate il giro in corso, fino a chi siede alla destra di chi ha iniziato la partita; poi giocate \textbf{ancora un giro completo}. In questi turni si possono chiudere altri set e prendere i token rimasti.

\begin{nota}{Esempio}
Chi ha iniziato la partita chiude il quinto set. Il resto del tavolo termina il giro; poi ogni partecipante, compreso chi ha chiuso, gioca un ultimo turno.
\end{nota}

Quando il mazzo degli Ambiti finisce, rimescolate gli scarti. Se un mazzo finisce e non si può ricostituire, la fine della partita scatta allo stesso modo.

\titolo{Vittoria}

Ogni partecipante somma i punti dei propri set e quelli dei token. Vince chi ha il punteggio più alto. In caso di parità, tra chi è a pari punti vince \textbf{la prima voce che urla «Lavoro!»}.

\clearpage
\titolo{Casi particolari}

\begin{itemize}
  \item Chi non può fare nessuna azione passa il turno.
  \item Si possono aprire quante pile si vuole, e migliorare una pila anche dopo averci messo la Persona.
  \item Il potere \potereQuattro{} piazza la seconda Persona su una pila che ne è ancora senza; quella Persona deve comunque avere coperte le altre due ambizioni.
  \item Una Persona messa in gioco da \potereQuattro{} non attiva il suo potere.
\end{itemize}

\titolo{Consigli per la prima partita}

\begin{itemize}
  \item \textbf{Non partire col primo set che ti riesce.} Tenere quattro o cinque carte in mano e cercare un Lavoro che somigli alle ambizioni della Persona paga quasi sempre.
  \item \textbf{Ma non aspettare solo il set perfetto.} Chi rinuncia a un set intero per farne uno da 16 resta indietro, anche nei token.
  \item \textbf{Guarda il mercato.} Le Persone con Ricerca sono le più difficili da formare: per questo hanno il potere più utile.
  \item \textbf{Chi chiude la partita non prende punti in più:} conta quanto vale ogni set.
\end{itemize}

\clearpage
% ================================================================ IN SINTESI
\titolo{Il turno in sintesi}

\begin{nota}{Una sola azione}
\textbf{1 Apri} una formazione · \textbf{2 Migliora} una formazione · \textbf{3 Inserisci} una Persona e attiva il suo potere · \textbf{4 Completa} un set: punti, token, formazioni negli scarti · \textbf{5 Pesca} 2 Ambiti, poi il mercato scorre. Alla fine: al massimo 6 carte in mano.
\end{nota}

\begin{nota}{Punti}
Ambizioni che combaciano col Lavoro: nessuna 5, una 7, due 10, tre 16. Token: 3 set della stessa categoria 5 e 3; 4 categorie diverse 5 e 3.
\end{nota}

\begin{nota}{Poteri}
\textbf{\potereUno}: \breveUno\par
\textbf{\potereDue}: \breveDue\par
\textbf{\potereTre}: \breveTre\par
\textbf{\potereQuattro}: \breveQuattro
\end{nota}

\begin{nota}{Fine}
Al primo quinto set si chiude il giro, poi si gioca un giro in più. Parità: vince la prima voce che urla «Lavoro!».
\end{nota}

\clearpage
% ================================================================ RETRO: PRINT AND PLAY
\titolo{Print and play}

Questo print and play si stampa in una volta sola: \textbf{A4, a colori, fronte e retro sul lato lungo, al 100\%} (dimensioni effettive, non "adatta alla pagina").

\begin{itemize}
  \item \textbf{Fogli 1-7:} le 60 Persone. \textbf{Fogli 8-14:} i 60 Ambiti.
  \item \textbf{Foglio 15:} 4 carte riassunto e 4 token, stampati sui due lati.
  \item \textbf{Fogli 16-18:} questo regolamento.
\end{itemize}

\textbf{Prova prima un foglio.} Stampate solo le prime due pagine e guardate il foglio controluce: i dorsi devono cadere dietro le carte.

\sottotitolo{Taglio al laser}
Carte, carte riassunto e token hanno \textbf{2 mm di abbondanza} per lato: un piccolo scarto del laser non lascia bordi bianchi. Agli angoli di ogni foglio ci sono quattro \textbf{crocini di registro} (cerchio e croce) per allineare il taglio alla stampa. I file di taglio stanno nell'archivio dei disegni laser, cartella \texttt{collocamento-print-and-play}. Cartoncino stampato: 80\%, 1000 mm/min, una passata. Il foglio delle carte riassunto va tagliato col fronte in alto.

\sottotitolo{Il regolamento}
Piegate a metà i tre fogli del regolamento, uno dentro l'altro, con la copertina fuori. Una pinzatura sulla piega lo tiene insieme.

\vfill
\begin{nota}{Le immagini}
Le immagini di questo prototipo sono generate con l'intelligenza artificiale: servono a provare il gioco, non sono la sua veste definitiva. Per la versione finale vogliamo affidare grafica e illustrazioni, con un compenso, a chi fa design e illustrazione per professione.
\end{nota}

{\fontsize{7}{9}\selectfont Collocamento è un prototipo di gioco di carte sul mercato del lavoro. Regolamento rivisto, ottobre 2026.\par}

\end{document}
